\section{Methods: A Ladder of Test-Time Search}
\label{sec:method}

All strategies share one policy LLM $\policy$, one per-step prompt template
(contrastive: candidate CWE descriptions are shown side by side), and one
structured-output schema. The \emph{only} thing that varies across strategies is
how they search (Figure~\ref{fig:arch}). This is what makes the compute-matched
comparison (\S\ref{sec:experiments}) a clean measurement of search structure
rather than prompt engineering.

\begin{figure}[t]
\centering
\begin{tikzpicture}[
  execute at begin picture={
    \definecolor{rewardblue}{RGB}{31,119,180}
    \definecolor{mctsred}{RGB}{214,39,40}
    \definecolor{mctsgreen}{RGB}{44,160,44}
  },
  font=\footnotesize,
  >=Stealth,
  box/.style={
    rounded corners=2pt,
    draw=black!60,
    line width=0.45pt,
    align=center,
    inner sep=3.5pt,
    fill=black!6
  },
  smallbox/.style={
    rounded corners=1.5pt,
    draw=black!60,
    line width=0.35pt,
    align=center,
    inner sep=2pt,
    fill=white
  },
  arrow/.style={->, line width=0.45pt},
  reward/.style={
    rounded corners=2pt,
    draw=rewardblue,
    line width=0.55pt,
    align=center,
    inner sep=3pt,
    fill=rewardblue!10
  }
]

\node[box, text width=1.25cm, minimum height=0.78cm] (input) at (0,0)
  {code patch\\(diff + msg)};

\node[box, text width=2.05cm, minimum height=1.75cm] (policy) at (2.35,0)
  {};
\node[align=center] at ([yshift=0.57cm]policy.center) {policy LLM\\(frozen)};
\node[smallbox, text width=1.55cm] (rank) at ([yshift=0.08cm]policy.center)
  {rank children};
\node[smallbox, text width=1.55cm, below=1pt of rank] (prop)
  {propose step};
\node[smallbox, text width=1.55cm, below=1pt of prop] (commit)
  {commit leaf};

\node[box, text width=1.45cm, minimum height=0.76cm] (controller) at (4.85,0)
  {search\\controller};

\node[anchor=west, align=left, inner sep=1pt] (ladder) at (5.76,0.02) {%
  Greedy\\
  Self-cons. (N)\\
  Best-of-N\\
  Beam (w)\\
  {\color{mctsred}MCTS-Tax (I)}\\
  {\color{mctsgreen}MCTS-Reason (I)}%
};
\draw[black!35, line width=0.35pt] ([xshift=-0.12cm,yshift=0.43cm]ladder.west) --
  ([xshift=-0.12cm,yshift=-0.52cm]ladder.west);

\node[reward, text width=1.55cm, minimum height=0.70cm] (reward) at (4.85,-1.45)
  {reward $r\in[0,1]$};
\node[smallbox, draw=rewardblue, text width=0.92cm,
  below left=2pt and -1pt of reward.south] (self) {self-eval};
\node[smallbox, draw=rewardblue, text width=1.05cm,
  below right=2pt and -1pt of reward.south] (judge) {external\\judge};

\node[box, text width=1.32cm, minimum height=0.78cm] (output) at (7.85,0)
  {CWE path\\category $\to$ leaf};

\draw[arrow] (input) -- (policy);
\draw[arrow] (policy) -- (controller);
\draw[arrow] (controller) -- (output);
\draw[arrow] (controller) -- (reward);
\draw[arrow, dashed, draw=rewardblue]
  (reward.west) .. controls +(left:0.75cm) and +(down:0.85cm) .. (controller.south);

\end{tikzpicture}

\caption{Shared architecture. Every strategy drives the same frozen policy LLM
through the same three operations and differs only in the search controller it
plugs in; the reward signal (self-evaluation or external judge) is used only by
the reward-guided strategies (best-of-$N$, beam, MCTS).}
\label{fig:arch}
\end{figure}

\begin{table}[t]
\centering
\small
\begin{tabular}{l l l l}
\toprule
 & Strategy & Searches over & Knob \\
\midrule
M0a & Flat & --- (one call, all leaves) & --- \\
M0b & Greedy & --- (argmax per level) & --- \\
M1 & Self-consistency & i.i.d.\ samples, majority & $N$ \\
M2 & Best-of-$N$ & i.i.d.\ samples, reward-max & $N$ \\
M3 & Beam & taxonomy, width-limited & $w$ \\
M4 & $\mctstax$ & taxonomy, reward-guided & $I$ \\
M5 & $\mctsreas$ & reasoning steps, reward-guided & $I$ \\
\bottomrule
\end{tabular}
\caption{The ladder of strategies, ordered by increasing search structure. All
share one policy, prompt, and output schema; each strategy's knob is set by the
compute-matched protocol (\S\ref{sec:experiments}).}
\label{tab:ladder}
\end{table}

\subsection{No-search baselines}
\textbf{Flat (M0a):} a single call lists all leaf candidates and the model picks
one. \textbf{Greedy-hierarchical (M0b):} one call per level, taking the argmax
child; defines $\Acc_{\text{Greedy}}$ and the policy-strength proxy $\polacc$.

\subsection{Sampling strategies}
\textbf{Self-Consistency (M1):} sample $N$ independent paths at temperature
$T>0$, return the majority leaf. \textbf{Best-of-$N$ (M2):} sample $N$ paths,
return the one maximizing the reward $r$ (\S\ref{sec:reward}). M1 isolates the
value of sampling alone; M2 adds the value of a reward signal.

\subsection{Structured search}
\textbf{Beam (M3):} maintain the top-$w$ partial paths per level under $r$;
deterministic given $\policy$. \textbf{$\mctstax$ (M4):} MCTS over the taxonomy.

\begin{algorithm}[t]
\caption{$\mctstax$: MCTS over the CWE taxonomy}
\label{alg:mcts}
\begin{algorithmic}[1]
\Require code $c$, taxonomy $\mathcal{T}$, policy $\policy$, reward $r$, budget $I$ iterations
\State root state $s_0 \gets (r)$ \Comment{partial path}
\For{$i = 1, \dots, I$}
  \State \textbf{Select}: descend by $Q(s,a) + c_{\text{exp}} \, P(a)\sqrt{N(s)+1}/(1+N(s,a))$ \Comment{PUCT; $P$ from policy scores}
  \State \textbf{Expand}: $\policy$ scores the children of the selected node; scores become priors $P$
  \State \textbf{Evaluate}: $r(s')$ scores the (partial/complete) path via self-eval or judge
  \State \textbf{Backprop}: update $Q,N$ up the search tree
\EndFor
\State \Return path through most-visited children (robust child)
\end{algorithmic}
\end{algorithm}

\textbf{$\mctsreas$ (M5):} the same selection/backup core, but search-tree nodes
are \emph{reasoning states} (a sequence of short free-form reasoning steps about
the patch), not CWE nodes. Each move either appends one reasoning step or
\emph{commits} to a single leaf CWE from the flat leaf list (prompts in
Appendix~\ref{app:prompts}); rollouts force a commit and score the resulting
root-to-leaf path, and the committed leaf is deterministically mapped to a
taxonomy path. Crucially, M5 never commits level-by-level: the coarse category is
implied by the final leaf, not chosen first. M4 vs.\ M5 is our test of whether
search over the output structure or over reasoning is more robust to early-level
compounding (Proposition~\ref{prop:cascade}).

\begin{figure}[t]
\centering
\begin{tikzpicture}[
  execute at begin picture={
    \definecolor{mctsred}{RGB}{214,39,40}
    \definecolor{mctsgreen}{RGB}{44,160,44}
  },
  font=\footnotesize,
  >=Stealth,
  node/.style={inner sep=1pt},
  tax/.style={
    circle,
    draw=black!60,
    line width=0.4pt,
    minimum size=0.34cm,
    inner sep=1pt,
    fill=black!5
  },
  leaf/.style={
    rounded corners=1.5pt,
    draw=black!60,
    line width=0.35pt,
    align=center,
    inner sep=2pt,
    fill=black!4
  },
  step/.style={
    rounded corners=2pt,
    draw=black!60,
    line width=0.4pt,
    align=center,
    inner sep=2.2pt,
    fill=black!5,
    text width=1.42cm
  },
  commit/.style={
    rounded corners=2pt,
    draw=mctsgreen,
    line width=0.45pt,
    align=center,
    inner sep=2.2pt,
    fill=mctsgreen!10,
    text width=1.10cm
  },
  edge/.style={->, line width=0.4pt},
  sel/.style={->, line width=0.85pt, draw=mctsred},
  faint/.style={black!28},
  greenedge/.style={draw=mctsgreen}
]

\node[align=center] at (1.85,1.15) {MCTS-Tax: search over taxonomy};
\node[tax] (root) at (1.85,0.55) {R};
\node[tax] (cat1) at (0.65,-0.15) {A};
\node[tax] (cat2) at (1.85,-0.15) {B};
\node[tax] (cat3) at (3.05,-0.15) {C};
\node[leaf] (a1) at (0.30,-0.90) {leaf};
\node[leaf] (a2) at (1.00,-0.90) {leaf};
\node[leaf, draw=black!28, text=black!38, fill=black!2] (b1) at (1.50,-0.90) {leaf};
\node[leaf, draw=black!28, text=black!38, fill=black!2] (b2) at (2.20,-0.90) {leaf};
\node[leaf] (c1) at (3.05,-0.90) {leaf};

\draw[sel] (root) -- (cat1);
\draw[edge] (root) -- (cat2);
\draw[edge] (root) -- (cat3);
\draw[edge] (cat1) -- (a1);
\draw[edge] (cat1) -- (a2);
\draw[edge, faint] (cat2) -- (b1);
\draw[edge, faint] (cat2) -- (b2);
\draw[edge] (cat3) -- (c1);

\node[align=center, text width=1.85cm, text=mctsred]
  at (1.00,-1.65) {commits at level 0\\unrecoverable};
\draw[mctsgreen, line width=0.55pt]
  ([xshift=-0.13cm,yshift=0.16cm]b2.north east) -- ++(0.08cm,-0.09cm) -- ++(0.17cm,0.20cm);
\node[align=center, text=black!45] at (1.85,-1.20) {unreachable};

\node[align=center] at (6.25,1.15) {MCTS-Reason: search over reasoning};
\node[step] (s0) at (4.65,0.45) {step:\\UAF risk};
\node[step] (s1) at (6.05,0.82) {step:\\realloc moves};
\node[step] (s2) at (6.05,0.06) {step:\\alias stale};
\node[commit] (k1) at (7.45,0.82) {commit:\\CWE-416};
\node[commit] (k2) at (7.45,0.06) {commit:\\CWE-416};
\node[leaf] (l1) at (8.25,0.82) {leaf};
\node[leaf] (l2) at (8.25,0.06) {leaf};

\draw[edge] (s0) -- (s1);
\draw[edge] (s0) -- (s2);
\draw[edge, greenedge] (s1) -- (k1);
\draw[edge, greenedge] (s2) -- (k2);
\draw[edge, greenedge] (k1) -- (l1);
\draw[edge, greenedge] (k2) -- (l2);

\node[align=center, text width=2.25cm, text=mctsgreen]
  at (6.35,-0.95) {defers commitment,\\picks leaf directly};

\end{tikzpicture}

\caption{The two MCTS search spaces. $\mctstax$ (left) walks the label tree and
must pick a coarse category before it has reasoned about the patch; a wrong
first move makes the gold leaf unreachable. $\mctsreas$ (right) searches over
short free-form reasoning steps and commits to a leaf directly, so the category
is implied by the final choice rather than locked in first.}
\label{fig:spaces}
\end{figure}

\subsection{Reward signal}
\label{sec:reward}
We deliberately use \emph{only} zero-shot reward sources, with no trained process
reward model (PRM), to keep the study GPU-free and reproducible:
\begin{itemize}
\item \textbf{Self-eval:} the policy $\policy$ scores its own (partial) path in
$[0,1]$.
\item \textbf{External judge:} a separate, typically stronger model scores the
path given $c$ and candidate CWE descriptions. We record the pair
$(\policy, \text{judge})$.
\end{itemize}
Comparing self-eval vs.\ judge directly probes whether the bottleneck is the
policy or the reward (RQ3/H3). A trained PRM is left to future work
(\S\ref{sec:conclusion}).

\subsection{Structured outputs}
Each policy call returns JSON --- per-candidate scores plus a one-sentence
rationale that is carried into the next level's prompt --- parsed defensively
with a small floor score for missing candidates, so malformed outputs never
crash the search (Appendix~\ref{app:prompts}). Conformal calibration of
path-level prediction sets from the score vectors is part of the final protocol.
